\subsection{BÀI TẬP TRẮC NGHIỆM}

\ind{PHẦN I.} \inden{Câu trắc nghiệm nhiều phương án lựa chọn. Mỗi câu hỏi học sinh chỉ chọn một phương án.}\\
\setcounter{ex}{0}
\Opensolutionfile{ans}[ans/1H7-B5-1]
\begin{ex}%[1H3Y5-1]
	Khẳng định nào sau đây là đúng?
	\choice
	{Khoảng cách giữa một đường thẳng và một mặt phẳng song song với nhau là khoảng cách từ một điểm thuộc mặt phẳng tới đường thẳng}
	{Khoảng cách từ một điểm tới một đường thẳng là khoảng cách từ điểm đó tới một điểm bất kì của đường thẳng}
	{Khoảng cách từ một điểm tới một mặt phẳng là khoảng cách từ điểm đó tới một điểm bất kì của mặt phẳng}
	{\True Khoảng cách giữa một đường thẳng và một mặt phẳng song song với nhau là khoảng cách từ một điểm thuộc đường thẳng tới mặt phẳng}
	\loigiai{
		Khoảng cách giữa một đường thẳng và một mặt phẳng song song với nhau là khoảng cách từ một điểm bất kì thuộc đường thẳng tới mặt phẳng đó.
	}
\end{ex}

\begin{ex}%[TT lần 1, 2017 - 2018, Chuyên Quốc Học Huế]%[Vũ Nguyễn Hoàng Anh, 12EX6]%[1H3Y5-1]%
	Trong các mệnh đề sau, mệnh đề nào \textbf{sai}?
	\choice
	{Khoảng cách giữa hai đường thẳng chéo nhau bằng khoảng cách giữa đường thẳng và mặt phẳng song song với nó đồng thời chứa đường thẳng kia}
	{\True Khoảng cách giữa hai đường thẳng chéo nhau bằng khoảng cách từ một điểm bất kì thuộc đường thẳng này đến đường thẳng kia}
	{Khoảng cách giữa hai đường thẳng chéo nhau bằng khoảng cách giữa hai mặt phẳng song song lần lượt chứa hai đường thẳng đó}
	{Khoảng cách giữa hai đường thẳng chéo nhau là độ dài đoạn vuông góc chung của hai đường thẳng đó}
	\loigiai{
		Mệnh đề \lq\lq Khoảng cách giữa hai đường thẳng chéo nhau bằng khoảng cách từ một điểm bất kì thuộc đường thẳng này đến đường thẳng kia\rq\rq\ là sai. Thật vậy, xét hình lập phương $ABCD.A'B'C'D'$ cạnh $1$. Theo định nghĩa $d(AB, A'D')=AA'=1$, tuy nhiên $d(B,A'D')=BA'=\sqrt{2}$.
	}
\end{ex}

\begin{ex}%[1H3Y5-4]
	Cho hình lập phương $ABCD.A'B'C'D'$ có cạnh bằng $a$. Khoảng cách giữa hai đường thẳng $BD$ và $A'C'$ bằng
	\choice
	{\True $a$}
	{$a\sqrt{2}$}
	{$\dfrac{a\sqrt{3}}{2}$}
	{$a\sqrt{3}$}
	\loigiai{
		\immini{
			Ta có
			$\mathrm{d}(BD,A'C')=\mathrm{d}((ABCD),(A'B'C'D'))=AA'=a$.
		}{
			\begin{tikzpicture}[line cap=round, line join=round,font=\footnotesize,>=stealth, scale=1]
				\tikzset{label style/.style={font=\footnotesize}}
				\tkzDefPoints{0/0/D',2.5/0/C',1/1/A'}
				\coordinate (B') at ($(A')+(C')-(D')$);
				\tkzDefSquare(D',C')    \tkzGetPoints{C}{D}
				\tkzDefSquare(A',B')    \tkzGetPoints{B}{A}
				\tkzDrawPolygon(A,B,B',C',D',D)
				\tkzDrawSegments(C,B C,D C,C' B,D)
				\tkzDrawSegments[dashed](A',A A',B' A',D' A',C')
				\tkzDrawPoints[fill=black](A',B',C',D',C,D,B,A)
				\tkzLabelPoints[above](A,B)
				\tkzLabelPoints[below](D',C')
				\tkzLabelPoints[left](A',D)
				\tkzLabelPoints[right](C,B')
			\end{tikzpicture}
		}
	}
\end{ex}

\begin{ex}%[1H3Y5-4]
	Cho hình lập phương $ABCD.A'B'C'D'$ có cạnh bằng $2a$. Tính khoảng cách giữa hai đường thẳng $AC$ và $A'B'$.
	\choice
	{$\dfrac{a\sqrt{3}}{2}$}
	{$\dfrac{a\sqrt{2}}{2}$}
	{$a\sqrt{2}$}
	{\True $2a$}
	\loigiai
	{
		\immini
		{
			Vì $A'B'\parallel AB$ nên $A'B'\parallel((ABCD))$.\\
			Vì vậy
			$$\mathrm{d}(A'B',AC)=\mathrm{d}(A'B',(ABCD))=\mathrm{d}(A',(ABCD))=AA'=2a.$$
		}
		{
			\begin{tikzpicture}[scale=0.7]
				\tkzDefPoints{0/0/A', 3/0/B', 2/-1.5/C'}
				\tkzDefPointBy[translation=from B' to C'](A')\tkzGetPoint{D'}
				\tkzCalcLength[cm](A',B')\tkzGetLength{rAB}
				\tkzDefShiftPoint[A'](90:\rAB){A}
				\tkzDefPointBy[translation=from A' to A](B')\tkzGetPoint{B}
				\tkzDefPointBy[translation=from A' to A](C')\tkzGetPoint{C}
				\tkzDefPointBy[translation=from A' to A](D')\tkzGetPoint{D}
				\tkzDrawPoints[fill=black](A,B,C,D,A',B',C',D')
				\tkzDrawPolygon(A,B,C,D)
				\tkzDrawSegments(B,B' C,C' D,D' B',C' C',D' A,C)
				\tkzDrawSegments[dashed](A,A' A',B' A',D')
				\tkzLabelPoints[above](A,B)
				\tkzLabelPoints[below](C',D',A')
				\tkzLabelPoints[left](D)
				\tkzLabelPoints[right](B')
				\tkzLabelPoints[below right](C)
				\tkzLabelSegment[pos=0.5, right](B,B'){$2a$}
			\end{tikzpicture}
		}
	}
\end{ex}

\begin{ex}%[1H3Y5-3]
	\immini[thm]
	{
		Cho hình hộp chữ nhật $ABCD.A'B'C'D'$ có ba kích thước $5$; $7$; $12$ (xem hình vẽ bên). Khoảng cách giữa hai mặt phẳng $(ABCD)$ và $(A'B'C'D')$ bằng
		\choice
		{$7$}
		{\True $12$}
		{$\sqrt{74}$}
		{$5$}
	}
	{
		\begin{tikzpicture}[line join=round,line cap=round,font=\footnotesize,scale=1]
			\coordinate[label=below left:$A$] (A) at (0,0);
			\coordinate[label=above left:$D$] (D) at (1,.8);
			\coordinate[label=below right:$B$] (B) at (3,0);
			\coordinate[label=above right:$C$] (C) at ($(B)-(A)+(D)$);
			\coordinate[label=above left:$D'$] (D') at ($(D)+(90:2.5)$);
			\coordinate[label=left:$A'$] (A') at ($(A)-(D)+(D')$);
			\coordinate[label=below right:$B'$] (B') at ($(B)-(D)+(D')$);
			\coordinate[label=above right:$C'$] (C') at ($(C)-(D)+(D')$);
			\draw (A')--(A)--(B)--(C)--(C')--(D')--(A')--(B')--(C') (B)--(B');
			\draw[dashed] (D)--(C) (D')--(D)--(A);
			\node at ($(A)!0.5!(D)$)[left]{$7$};
			\node at ($(A)!0.5!(B)$)[below]{$5$};
			\node at ($(A)!0.5!(A')$)[left]{$12$};
			\foreach \diem in {D,A,B,C,D',A',B',C'} \fill (\diem) circle(1pt);
		\end{tikzpicture}
	}
	\loigiai
	{
		Ta có khoảng cách giữa hai mặt phẳng $(ABCD)$ và $(A'B'C'D')$ là $AA'=12$.
	}
\end{ex}

\begin{ex}%[1H3B5-4]
	\immini[thm]
	{Cho hình hình lập phương $ABCD.A'B'C'D'$ có cạnh bằng $a$ như hình vẽ bên. Khoảng cách giữa hai đường thẳng $BD$ và $A'C'$ bằng
		\choice[2]
		{$\dfrac{\sqrt{3}a}{2}$}
		{$\sqrt{2}a$}
		{\True $a$}
		{$\sqrt{3}a$}
	}
	{\begin{tikzpicture}[scale=0.5, font=\footnotesize, line join=round, line cap=round, >=stealth]
			\tkzDefPoints{0/0/A, 1.5/-2/B, 6/0/D,7.5/-2/C}
			\coordinate (B') at ($(B)+(0,6)$);
			\coordinate (A') at ($(A)+(0,6)$);
			\coordinate (C') at ($(C)+(0,6)$);
			\coordinate (D') at ($(D)+(0,6)$);
			\tkzDefMidPoint(B,D)\tkzGetPoint{O}
			\tkzDefMidPoint(A',D)\tkzGetPoint{R}
			\tkzDefMidPoint(A',B)\tkzGetPoint{S}
			\tkzDefMidPoint(B',C)\tkzGetPoint{P}
			\tkzDefMidPoint(B',D')\tkzGetPoint{O'}
			\tkzDefMidPoint(C,D')\tkzGetPoint{Q}
			\tkzLabelPoints[left](A,A',B')
			\tkzLabelPoints[below](B)
			\tkzLabelPoints[right](C,C',D,D')
			\tkzDrawPoints[size=2](A, B, C,D, A',B',C',D')
			\tkzDrawSegments(A,B B',A' A',A B',B A',D' B',C' D',C' B,C C,C' A',C')
			\tkzDrawSegments[dashed](A,D C,D D,D' B,D)
		\end{tikzpicture}
	}
	\loigiai
	{Do giả thiết ta có $\left(ABCD\right)\parallel\left(A'B'C'D'\right)$ nên
		\allowdisplaybreaks
		\begin{eqnarray*}
			\mathrm{d}\left(BD; A'C'\right) &=& \mathrm{d}\left[\left(ABCD\right); \left(A'B'C'D'\right)\right]\\
			&=& \mathrm{d}\left(A; \left(A'B'C'D'\right)\right) = AA' = a.
		\end{eqnarray*}
	}
\end{ex}

\begin{ex}%[1H3B5-4]
	\immini[thm]{Cho hình lập phương $ABCD.A'B'C'D'$ có cạnh bằng $a$ (tham khảo hình vẽ bên). Khoảng cách giữa hai đường thẳng $BB'$ và $A'C'$ bằng
		\choice
		{ \True $\dfrac{a\sqrt 2}{2}$}
		{$a\sqrt{2}$ }
		{ $a\sqrt{3}$}
		{ $a$}
	}
	{
		\begin{tikzpicture}[scale=1,>=stealth, font=\footnotesize, line join=round, line cap=round]
			\tkzDefPoints{0/0/A,-1.3/-1.1/B,2/-1.1/C}
			\coordinate (D) at ($(A)+(C)-(B)$);
			\coordinate (A') at ($(A)-(0,2.5)$);
			\tkzDefPointsBy[translation=from A to A'](B,C,D){B'}{C'}{D'}
			\tkzDrawSegments(A,B B,C C,D D,A B,B' C',B' C',C C',D' D,D')
			\tkzDrawSegments[dashed](A',B' A',D' A,A')
			\tkzDrawPoints[fill=black,size=4](A,B,D,C,A',B',C',D')
			\tkzLabelPoints[above](A,D)
			\tkzLabelPoints[below](B',C')
			\tkzLabelPoints[left](B)
			\tkzLabelPoints[shift={(-0.65,0.4)}](A')
			\tkzLabelPoints[right](C,D')
		\end{tikzpicture}
	}
	\loigiai{
		\immini{
			Ta có $B'O \perp A'C'$ và $BB' \perp B'O$ nên $OB'$ chính là khoảng cách giữa $BB'$ và $A'C'$.\\
			Vậy khoảng cách giữa $BB'$ và $A'C'$ bằng $\dfrac{a\sqrt{2}}{2}$.
		}
		{
			\begin{tikzpicture}[scale=1,>=stealth, font=\footnotesize, line join=round, line cap=round]
				\tkzDefPoints{0/0/A,-1.3/-1.1/B,2/-1.1/C}
				\coordinate (D) at ($(A)+(C)-(B)$);
				\coordinate (A') at ($(A)-(0,2.5)$);
				\tkzDefPointsBy[translation=from A to A'](B,C,D){B'}{C'}{D'}
				\tkzInterLL(A',C')(B',D')\tkzGetPoint{O}
				\tkzDrawSegments(A,B B,C C,D D,A B,B' C',B' C',C C',D' D,D')
				\tkzDrawSegments[dashed](A',B' A',D' A,A' A',C' B',D')
				\tkzDrawPoints[fill=black,size=4](A,B,D,C,A',B',C',D')
				\tkzLabelPoints[above](A,D)
				\tkzLabelPoints[below](B',C',O)
				\tkzLabelPoints[left](B)
				\tkzLabelPoints[shift={(-0.65,0.4)}](A')
				\tkzLabelPoints[right](C,D')
				\tkzMarkRightAngles[size=0.17](B',O,C')
				\tkzMarkRightAngles[size=0.17](B,B',O)
				\tkzDrawPoints[fill=black](O)
			\end{tikzpicture}
		}
	}
\end{ex}

\begin{ex}%[1H3B5-2]
	Cho hình chóp $A.BCD$ có $AC\perp (BCD)$ và $BCD$ là tam giác đều cạnh bằng $a$, biết $AC=a\sqrt{2}$. Khoảng cách từ $A$ đến đường thẳng $BD$ bằng
	\choice
	{\True $\dfrac{a\sqrt{11}}{2}$}
	{$\dfrac{2a\sqrt{3}}{3}$}
	{$\dfrac{4a\sqrt{5}}{3}$}
	{$\dfrac{3a\sqrt{2}}{2}$}
	\loigiai{
		\immini{Gọi $H$ là trung điểm $BD$, ta có $CH=\dfrac{a\sqrt{3}}{2}$.\\
			Ta có $AH=\sqrt{AC^2+CH^2}=\sqrt{(a\sqrt{2})^2+\left(\dfrac{a\sqrt{3}}{2}\right)^2}=\dfrac{a\sqrt{11}}{2}$.\\
			Từ $AC\perp (BCD)$ và $CB=CD=a$ suy ra $AB=AD$, do đó $AH\perp BD$. Bởi vậy, khoảng cách từ $A$ đến đường thẳng $BD$ bằng $AH=\dfrac{a\sqrt{11}}{2}$.
		}
		{\begin{tikzpicture}[scale=0.8, font=\footnotesize, line join=round, line cap=round, >=stealth]
				\tkzDefPoints{0/0/C, 4/-3/B, 7/0/D}
				\coordinate (A) at ($(C)+(0,5)$);
				\coordinate (H) at ($(B)!.5!(D)$);
				\tkzDrawSegments[dashed](C,D C,H)
				\tkzDrawPolygon(A,C,B)
				\tkzDrawSegments(A,D D,B A,H)
				\tkzLabelPoints[left](A,C)
				\tkzDrawPoints(A,C,B,H,D)
				\tkzLabelPoints[below](B,D,H)
				\tkzMarkRightAngle(A,C,B)
				\tkzMarkRightAngle(A,C,D)
				\tkzMarkRightAngle(C,H,D)
				\tkzMarkRightAngle(A,H,D)
		\end{tikzpicture}}
	}
\end{ex}

\begin{ex}%7.48.%[1H2B5-2]
	Cho lăng trụ tam giác đều $ABC\cdot A'B'C'$ có tất cả các cạnh bằng $a$. Khoảng cách từ điểm $A$ đến đường thẳng $BC'$ bằng
	\choice
	{$\dfrac{a \sqrt{7}}{2}$}
	{\True $\dfrac{a \sqrt{14}}{4}$}
	{$\dfrac{a \sqrt{7}}{4}$}
	{$\dfrac{a \sqrt{14}}{2}$}
	\loigiai{
		\immini{Ta có: $BC'=AC'=a\sqrt{2}$.\\
			Gọi $M$ trung điểm $AB$ thì $C'M=\sqrt{BC'^2-MB^2}=\sqrt{2a^2-\tfrac{a^2}{4}}=\dfrac{a\sqrt{7}}{2}.$\\
			$S_{ABC'}=\dfrac{1}{2}C'M\cdot AB=\dfrac{1}{2}\cdot\dfrac{a\sqrt{7}}{2}\cdot a=\dfrac{a^2\sqrt{7}}{4}.$\\
			$S_{ABC'}=\dfrac{1}{2}\mathrm{d}(A,BC')\cdot BC'\Rightarrow \mathrm{d}(A,BC')=\dfrac{2\cdot S_{ABC'}}{BC'}=\dfrac{a\sqrt{14}}{4}.$}
		{\begin{tikzpicture}[scale=0.7]
				\coordinate(A) at (0,0);
				\coordinate(C) at (6,0);
				\coordinate(B) at (2,-3);
				\coordinate(A') at ($(A)+(90:4)$);
				\coordinate(B') at ($(B)+(90:4)$);
				\coordinate(C') at ($(C)+(90:4)$);
				\draw(A')--(A)--(B)--(C)--(C')--(A')--(B')--(B)(B')--(C')(B)--(C');
				\draw[dashed](A)--(C);
				\foreach \x/\y in {A/left,B/left,C/right,A'/left,B'/left,C'/right}{
					\node at (\x)[\y]{$\x$};
					\fill[red] (\x) circle (1.5pt);}
		\end{tikzpicture}}
	}
\end{ex}

\begin{ex}%[1H3Y5-3]
	Cho hình chóp $S.ABC$ có đáy $ABC$ là tam giác đều cạnh $a$. Cạnh bên $SA=a\sqrt{3}$ và vuông góc với mặt đáy $(ABC)$. Tính khoảng cách $d$ từ $A$ đến mặt phẳng $(SBC)$
	\choice
	{$d=a$}
	{$d=\dfrac{a\sqrt{3}}{2}$}
	{\True $d=\dfrac{a\sqrt{15}}{5}$}
	{$d=\dfrac{a\sqrt{5}}{5}$}
	\loigiai{
		\immini{Kẻ $AM\perp BC$. Vì tam giác $ABC$ đều nên $M$ là trung điểm của $BC$.\\
			Gọi $K$ là hình chiếu của $A$ trên $SM$, suy ra $AK \perp SM.\hfill (1)$\\
			Ta có $\heva{&AM \perp BC\\&BC \perp SA} \Rightarrow BC \perp (SAM) \Rightarrow BC \perp AK.\hfill(2)$\\
			Từ $(1)$ và $(2)$, suy ra $AK \perp (SBC)$ nên $\mathrm{d}(A,(SBC))=AK$.\\
			Trong $\triangle SAM$, có $AK=\dfrac{SA\cdot AM}{\sqrt{SA^2+AM^2}}=\dfrac{a\sqrt{15}}{5}$. }
		{\begin{tikzpicture}[line join=round,line cap=round,line width=.6pt,font=\footnotesize,scale=1]
				\coordinate[label=left:$A$] (A) at (0,0);
				\coordinate[label=below left:$B$] (B) at (1,-1);
				\coordinate[label=right:$C$] (C) at (4,0);
				\coordinate[label=above left:$S$] (S) at ($(A)+(90:4)$);
				\coordinate[label=below right:$M$] (M) at ($(C)!1/2!(B)$); % Định nghĩa điểm M thỏa mãn \vt{AM}=1/2*\vt{AB}
				\coordinate[label=above right:$K$] (K) at ($(S)!(A)!(M)$); % Điểm H là hình chiếu của điểm C trên đường thẳng AB
				\draw (A)--(B)--(C)--(S)--cycle (S)--(B) (S)--(M);
				\draw[dashed] (A)--(C) (A)--(M) (A)--(K);
				\draw ($ (A)!5pt!(C)$)--($(A)!2!($($(A)!5pt!(C)$)!.5!($(A)!5pt!(S)$)$)$)--($(A)!5pt!(S)$);
				\fill (A)circle(1.5pt) (B)circle(1.5pt) (C)circle(1.5pt) (S)circle(1.5pt) (M)circle(1.5pt) (K)circle(1.5pt);
			\end{tikzpicture}
		}
	}
\end{ex}

\begin{ex}%[1H3B5-3]
	\immini[thm]{Cho hình chóp $S.ABC$ có $SA$ vuông góc với mặt phẳng $(ABC)$, $SA=a$ và tam giác $ABC$ đều cạnh $a$ (tham khảo hình bên). Tính theo $a$ khoảng cách từ điểm $A$ đến mặt phẳng $(SBC)$.
		\choice
		{$a$}
		{$\dfrac{\sqrt{3}a}{2}$}
		{$\dfrac{a\sqrt{2}}{2}$}
		{\True  $\dfrac{\sqrt{21}a}{7}$}
	}{\begin{tikzpicture}[>=stealth,line join=round,line cap=round,font=\footnotesize,scale=0.7]
			\path
			(0,0)coordinate (A)++(0:5)coordinate(C)++(210:4)coordinate(B)
			(A)++(90:4)coordinate(S)
			;
			\draw (A)--(B)--(C)--(S)--cycle (S)--(B);
			\draw[dashed] (A)--(C);
			\foreach \x/\g in{A/180,B/-90,C/0,S/90}
			\fill[black](\x)circle(1pt) ($(\x)+(\g:3mm)$)node{$\x$};
			\draw pic[draw,blue,angle radius=2mm] {right angle = S--A--B};
	\end{tikzpicture}}
	\loigiai{
		\immini{
			Gọi $M$ là trung điểm cạnh $BC$, suy ra $AM\perp BC$.\\
			Ta có $\heva{&BC\perp AM\\&BC\perp SA}\Rightarrow BC\perp (SAM)$.\\
			Dựng đường cao $AH$ trong tam giác $SAM$, khi đó\\
			$\heva{&AH\perp BC\\&AH\perp SM}\Rightarrow AH\perp (SBC)$.\\
			Do đó $\mathrm{d}(A,(SBC))=AH$.\\
			Xét $\triangle ABM$ có $AM=\sqrt{AB^2-BM^2}=\dfrac{a\sqrt{3}}{2}$.\\
			Xét $\triangle SAM$ có $AH=\dfrac{SA\cdot AM}{\sqrt{AM^2+SA^2}}=\dfrac{a\sqrt{21}}{7}$.\\
			Vậy $\mathrm{d}(A,(SBC))=AH=\dfrac{a\sqrt{21}}{7}$.
		}
		{
			\begin{tikzpicture}[>=stealth,line join=round,line cap=round,font=\footnotesize,scale=1]
				\path
				(0,0)coordinate (A)++(0:5)coordinate(C)++(210:4)coordinate(B)
				(A)++(90:4)coordinate(S)
				($(B)!.5!(C)$)coordinate(M)
				($(S)!.4!(M)$)coordinate(H)
				;
				\draw (A)--(B)--(C)--(S)--cycle (M)--(S)--(B);
				\draw[dashed] (M)--(A)--(C) (A)--(H);
				\foreach \x/\g in{A/180,B/-90,C/0,S/90,M/-45,H/30}
				\fill[black](\x)circle(1pt) ($(\x)+(\g:3mm)$)node{$\x$};
				\draw pic[draw,blue,angle radius=2mm] {right angle = A--H--M};
				\draw pic[draw,blue,angle radius=2mm] {right angle = S--A--M};
			\end{tikzpicture}
		}
	}
\end{ex}

\begin{ex}%[1H3K5-3]
	Cho hình chóp $S.ABCD$ có $ABCD$ là hình chữ nhật tâm $I,AB=a,BC=a \sqrt{3}$; $H$ là trung điểm của $AI$. Biết $SH$ vuông góc với đáy và tam giác $SAC$ vuông tại $S$. Tính khoảng cách $d$ từ điểm $A$ đến mặt phẳng $(SBD)$.
	\choice
	{\True $d= \dfrac{a \sqrt{15}}{5}$}
	{$d=a \sqrt{15}$}
	{$d= \dfrac{a \sqrt{15}}{15}$}
	{$d= \dfrac{3a \sqrt{15}}{5}$}
	\loigiai{
		Dựng $HE \perp BD$ tại $E \in BD$; $HK \perp SE$ tại $K \in SE$
		\immini{suy ra $HK \perp (SBD)$ do đó \\ $ \mathrm{d}(A,(SBD))=2 \mathrm{d}(H,(SBD))=2HK$. \\
			Do $AC=BD=2a$ nên $IA=IB=AB=a \Rightarrow \triangle IAB$ đều \\
			Suy ra $HE=HI \sin60^{\circ}= \dfrac{a}{2} \cdot \dfrac{\sqrt{3}}{2}= \dfrac{a \sqrt{3}}{4}$. \\
			Lại có $SA^2=AH \cdot AC= \dfrac{a}{2} \cdot 2a=a^2 \Rightarrow SA=a$; \\ $SH= \sqrt{SA^2-AH^2}= \sqrt{a^2- \dfrac{a^2}{4}}= \dfrac{a \sqrt{3}}{2}$. \\
			Do đó
			$HK= \dfrac{SH \cdot HE}{\sqrt{SH^2+HE^2}}= \dfrac{a \sqrt{15}}{10}$. \\
			Vậy $ \mathrm{d}(A,(SBD))= \dfrac{a \sqrt{15}}{5}$.}
		{\begin{tikzpicture}[line width=0.6pt, line cap=round, line join=round,scale=1.3]
				\coordinate[label=above left:{$A$}] (A) at (0,0);
				\coordinate[label=below left:{$D$}] (D) at (-1.5,-1.7);
				\coordinate[label=right:{$B$}] (B) at (2.6,0);
				\coordinate[label=below right:{$C$}] (C) at ($(B)+(D)-(A)$);
				\coordinate[label=below:{$I$}] (I) at ($(A)!1/2!(C)$);
				\coordinate[label=below left:{$H$}] (H) at ($(A)!1/2!(I)$);
				\coordinate[label={$S$}] (S) at ($(H)+(0,2.5)$);
				\coordinate[label=below right:{$E$}] (E) at ($(I)!0.27!(B)$);
				\coordinate[label=right:{$K$}] (K) at ($(S)!0.85!(E)$);
				\tkzMarkRightAngle[size=0.18](C,S,A)
				\tkzMarkRightAngle[size=0.12](A,D,C)
				\tkzMarkRightAngle[size=0.11](S,H,E)
				\tkzMarkRightAngle[size=0.1](A,E,I)
				\tkzMarkRightAngle[size=0.1](H,K,E)
				\draw (S)--(D)--(C)--(B)--(S)--(C);
				\draw[dashed] (D)--(A)--(S)--(H)--(E)--(S) (C)--(A)--(B)--(D) (H)--(K);
				\tkzDrawPoints[size=3,fill=black](I,A,B,C,D,S,H,E,K)
		\end{tikzpicture}}
	}
\end{ex}

\begin{ex}%[1H3Y5-3]
	Cho hình chóp $S.ABCD$ có đáy là hình vuông tâm $O$, $SA$ vuông góc với mặt đáy. Hỏi mệnh đề nào sau đây là sai?
	\choice
	{$\mathrm{d}( C,( SAB ) )=\mathrm{d}( C,( SAD ) )$}
	{\True $\mathrm{d}( A,( SBD ) )=\mathrm{d}( B,( SAC ) )$}
	{$\mathrm{d}( B,( SCD ) )=2\mathrm{d}( O,( SCD ) )$}
	{$\mathrm{d}( S,( ABCD ) )=SA$}
	\loigiai
	{
		\immini
		{
			Vì $BC=2OC$ nên $\mathrm{d}\left( B,\left( SCD \right) \right)=2\mathrm{d}\left( O,\left( SCD \right) \right)$.\\
			Ta có $\mathrm{d}\left( C,\left( SAB \right) \right)=CB;\mathrm{d}\left( C,\left( SAD \right) \right)=CD$ mà $CB=CD$ nên $\mathrm{d}\left( C,\left( SAB \right) \right)=\mathrm{d}\left( C,\left( SAD \right) \right)$.\\
			Vì $SA$ vuông góc với mặt đáy nên $\mathrm{d}\left( S,\left( ABCD \right) \right)=SA$.
			Do đó đáp án sai là $\mathrm{d}( A,( SBD ) )=\mathrm{d}( B,( SAC ) )$.
		}
		{
			\begin{tikzpicture}[scale=0.8, line join=round, line cap=round]
				\tkzDefPoints{0/0/A,-1.3/-1.6/B,2.5/-1.6/C}
				\coordinate (D) at ($(A)+(C)-(B)$);
				\coordinate (S) at ($(A)+(0,3)$);
				\tkzInterLL(A,C)(B,D)\tkzGetPoint{O}
				\tkzDrawPolygon(S,B,C,D)
				\tkzDrawSegments(S,C)
				\tkzDrawSegments[dashed](A,S A,B A,D A,C B,D)
				\tkzDrawPoints[fill=black,size=4](D,C,A,B,S,O)
				\tkzLabelPoints[above](S)
				\tkzLabelPoints[below](A,B,C,O)
				\tkzLabelPoints[right](D)
			\end{tikzpicture}
		}
	}
\end{ex}

\begin{ex}%[1H3B5-2]
	Cho hình chóp $S.ABCD$ có đáy $ABCD$ là hình vuông cạnh $a$, $SA\perp (ABCD)$. Tính khoảng cách từ điểm $B$ đến mặt phẳng $(SAC)$.
	\choice
	{\True $\dfrac{a\sqrt{2}}{2}$}
	{$\dfrac{a\sqrt{2}}{4}$}
	{$\dfrac{a\sqrt{2}}{3}$}
	{$\dfrac{a}{2}$}
	\loigiai{
		\immini{
			Gọi $O$ là tâm của hình vuông $ABCD$. Ta có $BO\perp AC$ và $BO\perp SA$ nên $BO\perp (SAC)$.\\
			Vậy $\mathrm{d}(B,(SAC))=BO=\dfrac{a\sqrt{2}}{2}$.
		}{
			\begin{tikzpicture}[line cap=round, line join=round,font=\footnotesize,>=stealth, scale=1, x=0.3cm, y=0.3cm]
				\tikzset{label style/.style={font=\footnotesize}}
				\tkzDefPoints{0/0/D,7/0/C,3/3/A}
				\coordinate (B) at ($(A)+(C)-(D)$);
				\coordinate (S) at ($(A)+(0,6)$);
				\tkzDefMidPoint(A,C)\tkzGetPoint{O}
				\tkzDrawPolygon(S,B,C,D)
				\tkzDrawSegments(S,C)
				\tkzDrawSegments[dashed](A,S A,B A,D A,C B,D)
				\tkzDrawPoints[fill=black](D,C,A,B,S,O)
				\tkzLabelPoints[above](S)
				\tkzLabelPoints[left](A,D)
				\tkzLabelPoints[right](B,C)
				\tkzLabelPoints[below](O)
			\end{tikzpicture}
		}
	}
\end{ex}

\begin{ex}%[1H3B5-3]
	\immini[thm]{
		Cho hình chóp $S.ABC$ có đáy $ABC$ là tam giác đều cạnh $2a$, tam giác $SAB$ đều và nằm trong mặt phẳng vuông góc với mặt phẳng đáy. Khoảng cách từ điểm $C$ đến mặt phẳng $(SAB)$.
		\choice[2]
		{$ a\sqrt{6}$}
		{\True $ a\sqrt{3}$}
		{$\dfrac{a\sqrt{3}}{2}$}
		{$2a\sqrt{3}$}
	}
	{
		\begin{tikzpicture}[scale=0.45, font=\footnotesize, line join=round, line cap=round, >=stealth]
			\tkzDefPoints{0/0/A,7/0/B,2/-3/C}
			\coordinate (H) at ($(A)!0.5!(B)$);
			\coordinate (S) at ($(H)+(0,5)$);
			\tkzDrawPolygon(S,A,C,B)
			\tkzDrawSegments(S,C)
			\tkzDrawSegments[dashed](A,B)
			\tkzDrawPoints[size=2](S,A,C,B)
			\tkzLabelPoints[above](S)
			\tkzLabelPoints[below](C)
			\tkzLabelPoints[left](A)
			\tkzLabelPoints[right](B)
		\end{tikzpicture}
	}
	\loigiai{
		\immini{
			Trong $(SAB)$ kẻ $SH\perp AB$,  vì $(SAB)\perp (ABC)\Rightarrow SH\perp (ABC)$.\\
			Do đó $\mathrm{d}(S,(ABC))=CH=\dfrac{2a\sqrt{3}}{2}=a\sqrt{3}$ (do tam giác $SAB$ đều cạnh $2a$).
		}{
			\begin{tikzpicture}[scale=0.35, font=\footnotesize, line join=round, line cap=round, >=stealth]
				\tkzDefPoints{0/0/A,7/0/B,2/-3/C}
				\coordinate (H) at ($(A)!0.5!(B)$);
				\coordinate (S) at ($(H)+(0,5)$);
				\tkzDrawPolygon(S,A,C,B)
				\tkzDrawSegments(S,C)
				\tkzDrawSegments[dashed](A,B S,H)
				\tkzDrawPoints[fill=black](S,A,C,B,H)
				\tkzLabelPoints[above](S)
				\tkzLabelPoints[below](C,H)
				\tkzLabelPoints[left](A)
				\tkzLabelPoints[right](B)
				\tkzMarkRightAngles[size=0.4](S,H,B)
			\end{tikzpicture}
		}
	}
\end{ex}



\begin{ex}%[1H3B5-4]%[Nguyễn Văn Vũ]
	Cho hình lăng trụ đều $ABC.A'B'C'$ có tất cả các cạnh bằng $a$. Tính theo $a$ khoảng cách giữa hai đường thẳng $AA'$ và $BC$.
	\choice
	{$\dfrac{a\sqrt{3}}{4}$}
	{$a$}
	{$\dfrac{a\sqrt{2}}{2}$}
	{\True $\dfrac{a\sqrt{3}}{2}$}
	\loigiai{
		\immini{
			Gọi $M$ là trung điểm của $BC$. \\
			Do $ABC$ là tam giác đều cạnh $a$ nên ta có $AM=\dfrac{a\sqrt{3}}{2}$ và $AM\bot BC$. \hfill$(1)$\\
			Mặt khác ta lại có $ABC.{A}'{B}'{C}'$ là lăng trụ đều nên $A{A}'\bot \left(ABC\right)$$\Rightarrow A{A}'\bot AM$. \hfill$(2)$\\
			Từ (1) và (2) ta có $AM$ là đoạn vuông góc chung của $A{A}'$ và $BC$.\\
			Vậy $d\left(A{A}',BC\right)=AM$$=\dfrac{a\sqrt{3}}{2}$.}
		{\begin{tikzpicture}[scale=.7]
				\tkzDefPoints{0/0/A, 1/-1.2/B, 5/0/C, 0/5/z}
				\coordinate (A') at ($(A)+(z)$);
				\coordinate (B') at ($(B)+(z)$);
				\coordinate (C') at ($(C)+(z)$);
				\tkzDefMidPoint(B,C)
				\tkzGetPoint{M}
				\tkzDrawPolygon(A',B',C')
				\tkzDrawSegments(A,B B,C A,A' B,B' C,C')
				\tkzDrawSegments[dashed](A,C A,M)
				\tkzLabelPoints[left](A)
				\tkzLabelPoints[above](A',C')
				\tkzLabelPoints[below right](B')
				\tkzLabelPoints[below](B,C,M)
				\tkzMarkRightAngle(A,M,C)
		\end{tikzpicture}}
	}
\end{ex}



\begin{ex}%[1H3B5-3]
	Cho hình chóp đều $S.ABCD$, có đáy $ABCD$ là hình vuông cạnh là $2a$, cạnh bên bằng $3a$. Khoảng cách từ $A$ đến mặt phẳng $(SCD)$ bằng
	\choice
	{$\dfrac{a\sqrt {14}}{3}$}
	{\True $\dfrac{a\sqrt {14}}{2}$}
	{$a\sqrt {14}$}
	{$\dfrac{a\sqrt {14}}{4}$}
	\loigiai{
		\immini{
			Gọi tâm của đáy là $O$, $M$ là trung điểm của $CD$.\\
			Trong $(SOM)$, kẻ $OK$ vuông góc với $SM$ tại $K$.\\
			Khi đó ta có $\mathrm{d}(A,(SCD))=2\mathrm{d}(O,(SCD))=2OK$.\\
			$SO=\sqrt{SD^2-OD^2}=\sqrt{9a^2-2a^2}=\sqrt{7a^2}=a\sqrt{7}$.\\
			$\dfrac{1}{OK^2}=\dfrac{1}{OS^2}+\dfrac{1}{OM^2} =\dfrac{1}{7a^2}+\dfrac{1}{a^2}=\dfrac{8}{7a^2}$.\\
			Suy ra $\mathrm{d}(A,(SCD))=2OK=\dfrac{a\sqrt{14}}{2}$.
		}
		{
			\begin{tikzpicture}[scale=0.8]
				\tkzInit[xmin=-5,xmax=5,ymin=-3,ymax=7]
				\tkzDefPoints{0/0/O,3/1/D,-1/1/A,0/5/S}
				\tkzDefPointBy[translation = from A to O](O)
				\tkzGetPoint{C}
				\tkzDefPointBy[translation = from D to O](O)
				\tkzGetPoint{B}
				\tkzDefMidPoint(C,D)\tkzGetPoint{M}
				\coordinate (K) at ($(S)!3/4!(M)$);
				\tkzDrawSegments(S,B S,D S,C C,B C,D S,M)
				\tkzDrawSegments[dashed](A,B A,D A,S O,M B,D A,C S,O O,K)
				\tkzLabelPoints[left](A,B)
				\tkzLabelPoints[below](O)
				\tkzLabelPoints[right](D,C,M,K)
				\tkzLabelPoints[above](S)
				\tkzMarkRightAngle(S,K,O)
				\tkzDrawPoints(S,A,B,C,D,M,O,K)
			\end{tikzpicture}
		}
	}
\end{ex}



\begin{ex}%[1H3B5-4]
	Cho hình chóp $S.ABC$ có đáy $ABC$ là tam giác vuông cân tại $A$, mặt bên $SBC$ là tam giác đều cạnh $a$ và mặt phẳng $\left( SBC \right)$ vuông góc với mặt đáy. Tính theo $a$ khoảng cách giữa hai đường thẳng $SA$ và $BC$.
	\choice
	{$\dfrac{a\sqrt{22}}{11}$}
	{\True $\dfrac{a\sqrt{3}}{4}$}
	{ $\dfrac{a\sqrt{11}}{22}$ }
	{$\dfrac{a\sqrt{4}}{3}$}
	\loigiai{
		\immini{Gọi $H$ là trung điểm của $BC \Rightarrow SH\perp BC$ và $AH\perp BC$, nên suy ra $BC\perp (SHA)$.\\
			Vì $\left( SBC \right)$ vuông góc với mặt đáy $(AB)$ nên suy ra $SH \perp (ABC)$. \\
			Gọi $K$ là hình chiếu vuông góc của $H$ lên $SB$. Từ $BC\perp (SHA)\Rightarrow HK\perp BC$. Do đó $HK$ là đoạn vuông góc chung của $SA$ và $BC$. \\
			Có $ SH=\dfrac{a\sqrt{3}}{2}$, $AH=\dfrac{a}{2}$;\\
			$HK^2=\dfrac{SH^2.AH^2}{SH^2+AH^2}=\dfrac{3a^2}{16}\Rightarrow HK=\dfrac{a\sqrt{3}}{4}$.
		}
		{ \begin{tikzpicture}[line cap=round,line join=round, >=stealth,scale=1]
				\def \xa{1}
				\def \xb{-1}
				\def \y{4}
				\def \z{3}
				\coordinate (B) at (0,0);
				\coordinate (A) at ($(B)+(\xa,\xb)$);
				\coordinate (C) at ($ (B)+(\y,0) $);
				\coordinate (M) at ($ (B)!0.5!(C) $);
				\coordinate (G) at ($ (B)!2/3!(M) $);
				\coordinate (S) at ($ (G)+(0,\z) $);
				\coordinate (H) at (1.7,0);
				\tkzDefMidPoint(S,A)\tkzGetPoint{K}
				\draw [dashed] (B)--(C) (S)--(H) (A)--(H) (K)--(H);
				\draw (S)--(B)--(A)--(S)--(C)--(A);
				\tkzDrawPoints(S,A,B,C)
				\tkzLabelPoints[above right](H)
				\tkzLabelPoints[below right](C)
				\tkzLabelPoints[above](S)
				\tkzLabelPoints[above left](B,K)
				\tkzLabelPoints[below left](A)
			\end{tikzpicture}
		}
	}
\end{ex}

\begin{ex}%[1H3B5-4]
	Cho tứ diện đều $ABCD$ có cạnh đáy bằng $2$. Tính khoảng cách giữa hai đường thẳng $AB$ và $CD$.
	\choice
	{\True $\sqrt{2}$}
	{$\sqrt{3}$}
	{$1$}
	{$\dfrac{\sqrt{3}}{2}$}
	\loigiai
	{
		\immini{Gọi $M$ là trung điểm của $DC$.\\
			Do $\triangle ACD$ và $\triangle BCD$ đều nên $\heva{&CD \perp AM\\&CD \perp BM}\Rightarrow CD \perp (ABM)$ tại $M$.\\
			Kẻ $MN \perp AB$ tại $N$ thì $MN=\mathrm{d}(AB,CD)$.\\
			Do $\triangle MAB$ cân tại $M$ nên $N$ là trung điểm của $AB$.\\
			Xét $\triangle ANM$ vuông tại $N$. Ta có $AM=2\cdot \dfrac{\sqrt{3}}{2}=\sqrt{3}$; $AN=\dfrac{AB}{2}=1$. Suy ra $MN = \sqrt{AM^2-AN^2} = \sqrt{2}$.\\
			Vậy $\mathrm{d}(AB,CD)=\sqrt{2}$.}
		{\hspace{0.5cm}\begin{tikzpicture}[scale=.8]
				\tkzInit[ymin=-3.5,ymax=4,xmin=-1,xmax=6]
				\tkzClip
				\tkzDefPoints{0/0/A, 3.5/3/B, 4/-3/D, 5/0/C}
				\tkzDefMidPoint(D,C)\tkzGetPoint{M}
				\tkzDefMidPoint(B,A)\tkzGetPoint{N}
				\tkzDrawSegments[dashed](A,C A,M M,N)
				\tkzDrawPolygon(D,B,C)
				\tkzDrawSegments(D,A A,B B,M)
				\tkzLabelPoints[left](A,N)
				\tkzLabelPoints[right](C, M)
				\tkzLabelPoints[below](D)
				\tkzLabelPoints[above](B)
				\tkzMarkRightAngles(D,M,A B,N,M)
				\tkzMarkSegments[mark=|](D,M M,C N,B N,A)
				\tkzDrawPoints[fill=black](A,B,C,D,N,M)
		\end{tikzpicture}}
	}
\end{ex}


\begin{ex}%[1H3B5-4]%
	Cho hình chóp $S.ABCD$ có đáy $ABCD$ là hình vuông với đường chéo $AC=2a$, $SA$ vuông góc với mặt phẳng $(ABCD)$. Khoảng cách giữa hai đường thẳng $SB$ và $CD$ là
	\choice
	{\True $a\sqrt{2}$}
	{$a\sqrt{3}$}
	{$\dfrac{a}{\sqrt{2}}$}
	{$\dfrac{a}{\sqrt{3}}$}
	\loigiai{
		\immini{Ta có $\heva{&DA\perp SA\\&DA\perp AB}\Rightarrow DA\perp(SAB)$.\\
			Mặt khác $\heva{&CD\not\subset(SAB)\\&CD\parallel AB}\Rightarrow CD\parallel (SAB)$.\\
			Từ đó suy ra khoảng cách giữa $SB$ và $CD$ bằng khoảng cách giữa $(SAB)$ và $CD$ và bằng $DA$.\\
			$ABCD$ là hình vuông với đường chéo $AC=2a$ suy ra $DA=\sqrt{2}a$.\\
			Khoảng cách giữa hai đường thẳng $SB$ và $CD$ là $a\sqrt{2}$.}
		{\begin{tikzpicture}[line join = round, line cap = round,>=stealth,scale=0.8,font=\footnotesize]
				\path
				(0,0) coordinate (A)
				(-1,-2) coordinate (B)
				(4,0) coordinate (D)
				($(B)+(D)-(A)$) coordinate (C)
				($(A)+(0,3.5)$) coordinate (S);
				\draw(C)--(B)--(S)--(C)--(D)--(S);
				\draw[dashed](B)--(A)(S)--(A)--(D);
				\foreach \p/\g in {A/170, S/90, B/-90, C/-45, D/0}{\node at ($(\p)+(\g:3mm)$) {$\p$};\draw[fill=black] (\p) circle (1pt);}
	\end{tikzpicture}}}
\end{ex}

\begin{ex}%[1H3K5-4]
	Cho hình chóp $S.ABCD$ có đáy $ABCD$ là hình vuông cạnh bằng $a$, $SA$ vuông góc với mặt phẳng đáy $(ABCD)$, góc giữa đường thẳng $SC$ và mặt phẳng $(ABCD)$ bằng $45^\circ$. Khoảng cách giữa hai đường thẳng $SB$ và $AC$ bằng
	\choice
	{$a\sqrt{11}$}
	{\True $\dfrac{a\sqrt{10}}{5}$}
	{$\dfrac{2a\sqrt{11}}{3}$}
	{$a\sqrt{3}$}
	\loigiai{
		Ta có $\heva{&SA \perp (ABCD)\\ & SC \cap (ABCD)=C\\ & \widehat{SC,(ABCD)}=45^\circ}$
		$\Rightarrow \widehat{SCA}=45^\circ \Rightarrow \bigtriangleup SAC$ vuông cân tại $A \Rightarrow SA=AC=a\sqrt{2}$.\\
		Dựng $Bx\parallel AC \Rightarrow d(AC,SB)=d(AC,(SBx))$.\\
		Dựng $AE\perp Bx$, $AF\perp SE\ (1)$.\\
		Ta có $\heva{&Bx\perp AE\\ & Bx\perp SA}$
		$\Rightarrow Bx\perp (SAE)\Rightarrow Bx\perp AF \ (2)$.\\
		Từ $(1)$ và $(2)$ suy ra  $AF\perp (SBE)$\\
		$\Rightarrow \mathrm{d}(AC,(SBx))=\mathrm{d}(A,(SBx))=AF=\mathrm{d}(SB,AC)$\\
		Ta có $BE\parallel AC \Rightarrow BE \perp BD$.\\
		Lại có $\heva{& BO \perp AO\\ & AE\perp BE}$\\
		$\Rightarrow AEBO$ là hình chữ nhật $\Rightarrow AE=OB=\dfrac{a\sqrt{2}}{2}$\\
		$\bigtriangleup SAE$ vuông tại $A$ có $AF$ đường cao\\
		$\Rightarrow \dfrac{1}{AF^2}=\dfrac{1}{AE^2}+\dfrac{1}{SA^2}\Leftrightarrow AF=\dfrac{AE\times SA}{\sqrt{AE^2+SA^2}}=\dfrac{\dfrac{a\sqrt{2}}{2}\times a\sqrt{2}}{\sqrt{\left(\dfrac{a\sqrt{2}}{2}\right)^2+(a\sqrt{2})^2}}=\dfrac{a\sqrt{10}}{5}$\\
		Vậy $\mathrm{d}(SB,AC)=\dfrac{a\sqrt{10}}{5}$.
		\begin{center}
			\begin{tikzpicture}[line join =round, line cap = round,>= stealth,font=\footnotesize,line width=1pt,scale=1.2]
				\coordinate [label=above right:$A$] (A) at (0,0);
				\coordinate [label=below:$B$] (B) at (-1,-2);
				\coordinate [label=right:$D$] (D) at (7,0);
				\coordinate [label=right:$C$] (C) at (6,-2);
				\coordinate [label=above:$O$] (O) at (3,-1);
				\coordinate[label=left:$S$] (S) at (0,3);
				\coordinate[label= below left:$E$] (E) at (-1.99,-1.19);
				\coordinate[label=right:$F$] (F) at ($(S)!0.4!0:(E)$);
				\coordinate[label=left:$x$] (x) at (-2.22,-1);
				\draw[dashed] (S)--(A)--(D)--(B)--(A)--(C);
				\draw (S)--(B)--(C)--(S)--(D)--(C);
				\draw (B)--(x) (S)--(E);
				\draw[dashed] (A)--(E) (A)--(F) ;
				\tkzMarkRightAngle (A,F,E);
				\tkzMarkRightAngle (A,E,B);
			\end{tikzpicture}
		\end{center}
	}
\end{ex}

\begin{ex}%[Thi thử, Triệu Quang Phục - Hưng Yên Lần 2, 2019]%[Lê Quốc Hiệp, 12EX5-2019]%[1H3K5-4]%
	Cho tứ diện $OABC$ có $OA$, $OB$, $OC$ đôi một vuông góc với nhau, $OA=a,~OB=OC=2a$.  Gọi $M$ là trung điểm của $BC$. Khoảng cách giữa hai đường thẳng $OM$ và $AB$ bằng
	\choice
	{$\dfrac{a\sqrt{2}}{2}$}
	{$\dfrac{2a\sqrt{5}}{5}$}
	{$a$}
	{\True $\dfrac{a\sqrt{6}}{3}$}
	\loigiai
	{
		\immini
		{
			Qua $B$ kẻ đường thẳng song song $OM$ cắt $OC$ tại $D$, từ đó suy ra $OM\parallel(ABD)$.\\
			Khi đó
			\[\mathrm{d}(OM,AB)=\mathrm{d}(OM,(ABD))=\mathrm{d}(O,(ABD)).\]
			Gọi $H$ là hình chiếu vuông góc của $O$ lên $BD$ và $K$ lần lượt là hình chiếu vuông góc của $O$ lên $AH$.\\
		}
		{\begin{tikzpicture}[line cap=round,line join=round,font=\footnotesize,>=stealth,scale=0.75]
				\fill (0,0) coordinate [label=below:$O$] (O) circle(1pt)
				(0:4) coordinate [label=below right:$C$] (C) circle(1pt)
				(90:4) coordinate [label=above:$A$] (A) circle(1pt)
				(-40:3) coordinate [label=below right:$B$] (B) circle(1pt)
				($(B)!0.5!(C)$) coordinate [label=below right:$M$] (M) circle(1pt);
				\tkzDefLine[parallel=through B](O,M)\tkzGetPoint{b}
				\tkzInterLL(B,b)(O,C)\tkzGetPoint{D}
				\tkzDefLine[parallel=through O](B,C)\tkzGetPoint{o}
				\tkzInterLL(O,o)(B,D)\tkzGetPoint{H}
				\fill($(A)!0.6!(H)$) coordinate [label=left:$K$] (K) circle(1pt);
				\fill (D) node[shift={(-120:2ex)}]{$D$} circle(1pt)
				(H) node[shift={(-120:2ex)}]{$H$} circle(1pt);
				\draw (H)--(A)--(D)--(B)--(C)--(A)--(B);
				\draw[dashed] (A)--(C)--(D) (A)--(O)--(M) (H)--(O)--(K) (B)--(O);
				\tkzMarkRightAngles[size=0.2](O,H,B O,K,H)
		\end{tikzpicture}}
		\noindent Ta có
		\[BD\perp OH\Rightarrow BD\perp(AOH)\Rightarrow BD\perp OK.\]
		Mà
		\[\heva{&OK\perp BD\\&OK\perp AH}\Rightarrow OK\perp (ABD)\Rightarrow\mathrm{d}(O,(ABD))=OK.\]
		Ta có
		\[OH=\dfrac{1}{2}BD=OM=\dfrac{OB\cdot OC}{BC}=\dfrac{2a\cdot2a}{2\sqrt{2}a}=a\sqrt{2}.\]
		Xét tam giác vuông $AOH$ ta có
		\[\dfrac{1}{OK^2}=\dfrac{1}{OH^2}+\dfrac{1}{OA^2}=\dfrac{1}{2a^2}+\dfrac{1}{a^2}=\dfrac{3}{2a^2}.\]
		Suy ra $OK=\dfrac{a\sqrt{6}}{3}$.
	}
\end{ex}

\begin{ex}%[1H3K5-4]
	Cho hình chóp $S.ABCD$  có đáy $ABCD$  là hình vuông tâm $O$  cạnh  $a$, $SO$  vuông góc với mặt phẳng $ABCD$  và $SO=a$.  Khoảng cách giữa $SC$  và $AB$  bằng
	\choice
	{\True $\dfrac{2a\sqrt{5}}{15}$}
	{$\dfrac{2a\sqrt{3}}{15}$}
	{$\dfrac{a\sqrt{3}}{15}$}
	{$\dfrac{a\sqrt{5}}{5}$}
	\loigiai
	{ \immini{ Vì $AB \parallel (SCD)$ nên
			$\mathrm{d}(AB,SC) = \mathrm{d}(AB,(SCD)) = \mathrm{d}(A,(SCD))=2\mathrm{d}(O,(SCD))$. \\ Gọi $K$ là trung điểm của $CD$. Khi đó $OK \perp CD \Rightarrow CD \perp (SOK)$. Gọi $H$ là hình chiếu vuông góc của $O$ lên $SK \Rightarrow OH \perp (SCD) \Rightarrow \mathrm{d}(O,(SCD)) = OH$. Ta có $\dfrac{1}{OH^2} =  \dfrac{1}{OK^2} + \dfrac{1}{OS^2} = \dfrac{1}{\left(\frac{a}{2}\right)^2}+\dfrac{1}{a^2}=\dfrac{5}{a^2} \Rightarrow OH = \dfrac{a\sqrt{5}}{5}$.\\
			Vậy $\mathrm{d}(AB,SC)  = 2 OH =  \dfrac{2a\sqrt{5}}{5}$.
		}{
			\begin{tikzpicture}[scale=.8]
				\tkzDefPoints{0/0/A, -2/-2/B, 3/-2/C}
				\coordinate (D) at ($(A)+(C)-(B)$);
				\coordinate (O) at ($(A)!.5!(C)$);
				\coordinate (S) at ($(O)+(0,5)$);
				\tkzDefMidPoint(D,C)\tkzGetPoint{K}
				\coordinate (H) at ($(S)!.7!(K)$);
				\tkzDrawSegments[dashed](S,A A,B A,D A,C B,D S,O O,K O,H)
				\tkzDrawPolygon(S,C,D)
				\tkzDrawSegments(S,B B,C S,K)
				\tkzLabelPoints[left](A,B)
				\tkzLabelPoints[right](C,D,K,H)
				\tkzLabelPoints[above](S)
				\tkzLabelPoints[below](O)
				\draw(S) circle (1pt) ;
				\draw(O) circle (1pt) ;
				\draw(H) circle (1pt) ;
				\draw(K) circle (1pt) ;
				\draw(A) circle (1pt) ;
				\draw(B) circle (1pt) ;
				\draw(C) circle (1pt) ;
				\draw(D) circle (1pt) ;
			\end{tikzpicture}
		}
	}
\end{ex}

\begin{ex}%[1H3K5-3]
	Cho hình chóp $S.ABCD$ có $SA\perp(ABCD)$, đáy $ABCD$ là hình thang vuông có chiều cao $AB=a$. Gọi $I$ là $J$ lần lượt là trung điểm của $AB$ và $CD$. Tính khoảng cách giữa $IJ$ và $(SAD)$.
	\choice
	{$\dfrac{a\sqrt{3}}{2}$}
	{$\dfrac{a}{3}$}
	{$\dfrac{a\sqrt{2}}{2}$}
	{\True $\dfrac{a}{2}$}
	\loigiai{
		\immini{Ta có $BC\parallel AD$ nên\\ $\mathrm{d}(IJ,(SAD))=d(I,(SAD))=AI=\dfrac{a}{2}$\\
			(Do $AI\perp(SAD)$).}{\begin{tikzpicture}[scale=1]
				\tkzDefPoints{0/0/A, -1/-1/B, 1/-1/C}
				\coordinate (E) at ($(A)+(C)-(B)$);
				\coordinate (S) at ($(A)+(0,1)$);
				\coordinate (D) at ($(A)!1.3!(E)$);
				\coordinate (I) at ($(A)!0.5!(B)$);
				\coordinate (J) at ($(C)!0.5!(D)$);
				\tkzDrawSegments[dashed](S,A A,B A,D I,J)
				\tkzDrawPolygon(S,C,D)
				\tkzDrawSegments(S,B B,C)
				\tkzDrawPoints(I,J)
				\tkzLabelPoints[left](A,B)
				\tkzLabelPoints[right](C,D,J)
				\tkzLabelPoints[above](S)
				\tkzLabelPoints[below](I)
				\tkzMarkRightAngle(D,A,S)
				\tkzMarkRightAngle(S,A,B)
				\tkzMarkRightAngle(A,B,C)
				\tkzMarkRightAngle(B,A,D)
	\end{tikzpicture}}}
\end{ex}

\begin{ex}%[1H3K5-3]
	Cho lăng trụ tứ giác đều $ABCD.A'B'C'D'$ có cạnh đáy bằng $a$. Gọi $M$, $N$, $P$ là trung điểm của các cạnh $AD$, $CD$, $A'D'$. Tính khoảng cách giữa $CC'$ và mặt phẳng $(MNP)$.
	\choice
	{$\dfrac{{a\sqrt 3 }}{3}$}
	{\True $\dfrac{{a\sqrt 2 }}{4}$}
	{$\dfrac{a}{{\sqrt 2 }}$}
	{$a\sqrt 2 $}
	\loigiai{
		\begin{center}
			\begin{tikzpicture}[scale = 0.7, font=\footnotesize, line join=round, line cap=round, >=stealth]
				\tkzDefPoints{0/0/A, -1.5/-1.5/B, 4/0/D, 0/4/A'};
				\coordinate (C) at ($ (B)+(D)-(A) $);
				\coordinate (C') at ($ (A')+(C)-(A) $);
				\coordinate (B') at ($ (B)+(C')-(C) $);
				\coordinate (D') at ($ (A')+(C')-(B') $);
				\tkzDrawSegments(C,B C,D C,C' A',B' B,B' B',C' C',D' A',D' D,D');
				\tkzDrawSegments[dashed](A,B A,D A,A');
				\tkzDrawPoints(A,B,C,D,A',B',C',D');
				\tkzLabelPoints[above](A',B',D');
				\tkzLabelPoints[below right](C');
				\tkzLabelPoints[below](A,B,C,D);
				\tkzDefMidPoint(A',D')\tkzGetPoint{P};
				\tkzDefMidPoint(A,D)\tkzGetPoint{M};
				\tkzDefMidPoint(C,D)\tkzGetPoint{N};
				\tkzDrawPoints(M,N,P);
				\tkzLabelPoints[below](M,N);
				\tkzLabelPoints[above](P);
				\tkzDrawSegments[dashed](P,M M,N N,P);
			\end{tikzpicture}
			\begin{tikzpicture}[scale = 0.7, font=\footnotesize, line join=round, line cap=round, >=stealth]
				\tkzDefPoints{0/0/B, 0/4/A, 4/4/D, 4/0/C, 2/4/M, 4/2/N, 6/0/I};
				\tkzDefMidPoint(N,I)\tkzGetPoint{H};
				\tkzDrawSegments(A,D A,B D,C B,C M,I C,H C,I);
				\tkzDrawPoints(A,B,C,D,M,N,I,H);
				\tkzLabelPoints[above](A,D,M,H);
				\tkzLabelPoints[below](B,C,I);
				\tkzLabelPoints[left](N);
				\tkzMarkRightAngle(C,H,I);
			\end{tikzpicture}
		\end{center}
		Ta có $\heva{ & MP\parallel CC' \\ & MP \subset \left( {MNP} \right) \\ & CC' \not\subset \left( {MNP} \right)}\Rightarrow CC'\parallel \left( {MNP} \right)$\\
		$ \Rightarrow \mathrm{d}\left( {CC';\left( {MNP} \right)} \right) = \mathrm{d}\left( {C;\left( {MNP} \right)} \right) = CH = \dfrac{a}{2}\cdot\dfrac{{\sqrt 2 }}{2} = \dfrac{{a\sqrt 2 }}{4}$.
	}
\end{ex}


\Closesolutionfile{ans}

\ind{PHẦN II.} \inden{Câu trắc nghiệm đúng sai. Trong mỗi ý a), b), c), d) ở mỗi câu, học sinh chọn đúng hoặc sai.}\\
\Opensolutionfile{ans}[ans/1H7-B5-2]
\begin{ex}%[1H8V5-5]
	\immini[thm]{Cho hình chóp $S.ABCD$ có đáy là hình thang, $AB=2a$, $AD=DC=CB=a$, $SA$ vuông góc với mặt phẳng đáy và $SA=3a$ (Hình bên). Gọi $M$ là trung điểm cạnh $AB$. Kẻ $AH\bot SC$. Các mệnh đề sau đúng hay sai?
		\choiceTF
		{$DM$ là đoạn vuông góc chung của $SB$ và $CD$}
		{\True Khoảng cách từ điểm $A$ đến mp $\left(SBC\right)$ bằng $\dfrac{3a}{2}$}
		{Khoảng cách $\mathrm{d}\left(M,\left(SBC\right) \right)=\dfrac{2}{3}\mathrm{d}\left(A,\left(SBC\right) \right)$}
		{\True Khoảng cách giữa cạnh $SB$ và $DM$ bằng $\dfrac{3a}{4}$}}{
		\begin{tikzpicture}[line join=round,line cap=round,>=stealth,font=\footnotesize,scale=.6]
			\coordinate (A) at (0,0);
			\coordinate (D) at (1.3,-2);
			\coordinate (B) at (7,0);
			\coordinate (C) at ($(D)+.4*(B)-(A)$);
			\coordinate (S) at ($(A)+(0,5)$);
			\coordinate (M) at ($(A)!.5!(B)$);
			\coordinate (H) at ($(S)!(A)!(C)$);
			\coordinate (O) at (intersection of A--C and B--D);
			\draw (S)--(A) (S)--(B) (S)--(C) (S)--(D) (B)--(C)--(D) (A)--(D);
			\draw[dashed] (A)--(B) (D)--(M) (A)--(H);
			\foreach \i/\g in {S/90,A/180,B/0,C/-90,D/-130,M/90,H/30}{\draw[fill=black](\i) circle (1.5pt) ($(\i)+(\g:4mm)$) node[scale=1]{$\i$};}
	\end{tikzpicture}}
	\loigiai{
		\begin{itemchoice}
			\itemch Vì $AM = MD = MC = MB$ nên $ADCB$ là nửa lục giác đều cạnh $a$, tâm $M$.\\
			$\Rightarrow \widehat{MDC} = {60}^\circ$.\\
			$\Rightarrow DM$ không vuông góc $CD$.\\
			$\Rightarrow DM$ không là đoạn vuông góc chung.
			\itemch Ta có $\heva{&AB \bot AC\\&CB'\bot SA} \Rightarrow CB \bot \left(SAC\right)$.\\
			$\Rightarrow CB \bot AH$ mà $AH \bot SC$\\
			$\Rightarrow AH \bot \left(SCB\right) \Rightarrow \mathrm{d}\left(A;\left(SBC\right)\right) = AH$.\\
			Ta lại có\\
			$\dfrac{1}{AH^2} = \dfrac{1}{SA^2} + \dfrac{1}{AC^2} = \dfrac{1}{\left(3a\right)^2} + \dfrac{1}{\left(2a\right)^2-a^2} = \dfrac{1}{2} \mathrm{d}\left(A;\left(SBC\right)\right) = \dfrac{3a}{4}$\\
			$\Rightarrow AH = \dfrac{3a}{2}$.
			\itemch Vì $M$ là trung điểm $AB$ nên $\dfrac{\mathrm{d}\left(M;\left(SBC\right)\right)}{\mathrm{d}\left(A;\left(SBD\right)\right)} = \dfrac{1}{2}$\\
			$\Rightarrow \mathrm{d}\left(M;\left(SBC\right)\right) = \dfrac{1}{2} \left(A;\left(SBC\right)\right)$
			\itemch Gọi $N$ là trung điểm $SA$ $\Rightarrow MN \parallel SB$\\
			$\Rightarrow \mathrm{d} \left(MD;SB\right) = \mathrm{d} \left(\left(MND\right);SB\right) = \mathrm{d} \left(B;\left(MND\right)\right) = \mathrm{d} \left(A;\left(DMN\right)\right)$
		\end{itemchoice}
	}
\end{ex}
%%%==============EX_2============%%%
\begin{ex}%[1H8V5-4]
	Cho hình chóp $S.ABCD$ có đáy $ABCD$ là hình vuông cạnh $a$. $SAD$ là tam giác đều và nằm trong mặt phẳng vuông góc với đáy. Gọi $M$ là trung điểm $AD$, $H$ là chân đường vuông góc hạ từ $M$ lên $SO$. Xét tính đúng sai của các mệnh đề sau:
	\choiceTF
	{\True Khoảng cách từ điểm $S$ đến mặt phẳng $\left(ABCD\right)$ là $SM$}
	{Khoảng cách từ điểm $D$ đến mặt phẳng $\left(SAC\right)$ bằng $2MH$}
	{\True Khoảng cách từ B đến mặt phẳng $\left(SAD\right)$ bằng $a$}
	{\True Khoảng cách giữa $AB$ và $SD$ là $\dfrac{a\sqrt{3}}{2}$}
	\loigiai{
		\begin{tikzpicture}[line join=round,line cap=round,>=stealth,font=\footnotesize,scale=.8]
			\coordinate (A) at (0,0);
			\coordinate (B) at (-3,-2.5);
			\coordinate (D) at (8,0);
			\coordinate (C) at ($(D)+(B)-(A)$);
			\coordinate (M) at ($(A)!.5!(D)$);
			\coordinate (S) at ($(M)+(0,7.33)$);
			\coordinate (O) at (intersection of A--C and B--D);
			\coordinate (H) at ($(S)!(M)!(O)$);
			\coordinate (N) at ($(A)!.5!(O)$);
			\coordinate (P) at ($(S)!(M)!(N)$);
			\coordinate (T) at ($(S)!(A)!(D)$);
			\draw(S)--(B) (S)--(C) (S)--(D) (B)--(C)--(D);
			\draw[dashed](S)--(A) (A)--(B) (A)--(D) (A)--(C) (S)--(M) (A)--(T)
			(D)--(B) (S)--(O) (M)--(H) (S)--(N) (M)--(N) (M)--(P);
			\pic[draw,angle radius=2mm] {right angle = S--M--D} 
			pic[draw,angle radius=2mm] {right angle = O--H--M}
			pic[draw,angle radius=2mm] {right angle = M--P--N}
			pic[draw,angle radius=2mm] {right angle = M--N--O}
			pic[draw,angle radius=2mm] {right angle = A--T--D}
			;
			\foreach \i/\g in {S/90,A/130,B/180,C/-45,D/0,M/-30,O/-90,N/-140,P/130,T/45}{\draw[fill=black](\i) circle (1.5pt) ($(\i)+(\g:4.5mm)$) node[scale=1]{$\i$};}
			\foreach \i/\g in {H/50}{\draw[fill=black](\i) circle (1.5pt) ($(\i)+(\g:6mm)$) node[scale=1]{$\i$};}
		\end{tikzpicture}
		\begin{itemchoice}
			\itemch $\heva{&\left(SAD\right) \cap \left(ABCD\right) = AD\\
				&\left(SAD\right) \bot \left(ABCD\right)\\
				& SM \subset \left(SAD\right), SM \bot AD
			}$ $\Rightarrow SM \bot \left(SAD\right)$\\
			$\Rightarrow \mathrm{d} \left(S;\left(ABCD\right)\right) = SM$.
			\itemch Gọi $N$ là trung điểm $AO$ $\Rightarrow MN \bot AO$.\\
			Kẻ $MP \bot SN$\\
			Ta có $\heva{&AO \bot MN\\&AO \bot SM}$ $\Rightarrow AO \bot \left(SMN\right)$\\
			$\Rightarrow AO \bot MP$ mà $MP \bot SN$\\
			$\Rightarrow MP\bot \left(SAC\right) \Rightarrow \mathrm{d}\left(M;\left(SAC\right)\right) = MP$.\\
			Ta có $\dfrac{\mathrm{d}\left(M;\left(SAC\right)\right)}{\mathrm{d}\left(D;\left(SAC\right)\right)} = \dfrac{1}{2} \Rightarrow \mathrm{d}\left(D;\left(SAC\right)\right) = 2 MP$.
			\itemch Vì $\heva{&AB \bot AD \\& AB \bot SM } \Rightarrow AB \bot \left(SAD\right) \Rightarrow AB = a = \mathrm{d}\left(B;\left(SAD\right)\right)$.
			\itemch Kẻ $AT \bot SD$.\\
			Vì $\heva{&AT \bot SD\\&AT\bot AB \left(\text{do } AB \bot \left(SAD\right) \text{ và } AT \subset \left(SAD\right) \right)}$\\
			$\Rightarrow SD = \mathrm{d}\left(SD;AB\right)=\dfrac{a\sqrt{3}}{2}$
		\end{itemchoice}
	}
\end{ex}
%%%==============EX_3============%%%
\begin{ex}%[1H8V5-3]
	Cho hình chóp $S.ABC$ có $SA\bot \left(ABC\right)$, $SA=a$, tam giác $ABC$ là tam giác đều cạnh $a$.
	\choiceTF
	{Khoảng cách giữa hai đường thẳng $SA$ và $BC$ là $\dfrac{a\sqrt{2}}{2}$}
	{Khoảng cách từ $C$ đến mặt phẳng $\left(SAB\right)$ là $\dfrac{a\sqrt{3}}{2}$}
	{Khoảng cách từ $A$ đến mặt phẳng $\left(SBC\right)$ là $\dfrac{a\sqrt{21}}{3}$}
	{Gọi $G$ là trọng tâm tam giác $ABC$. Khoảng cách từ $G$ đến mặt phẳng $\left(SBC\right)$ là $\dfrac{a\sqrt{21}}{21}$}
	\loigiai{
		\begin{center}
			\begin{tikzpicture}[line join=round,line cap=round,>=stealth,font=\footnotesize,scale=.7]
				\coordinate (A) at (-2,0);
				\coordinate (B) at (-.5,-2);
				\coordinate (C) at (3,0);
				\coordinate (S) at ($(A)+(0,4)$);
				\coordinate (M) at ($(B)!.5!(C)$);
				\coordinate (N) at ($(A)!.5!(B)$);
				\coordinate (H) at ($(S)!(A)!(M)$);
				\coordinate (G) at (intersection of N--C and A--M);
				\draw (S)--(A) (A)--(B) (S)--(C) (B)--(C) (S)--(B) (M)--(S);
				\draw[dashed] (A)--(C) (A)--(M) (A)--(H) (C)--(N);
				\pic[draw,angle radius=2mm] {right angle = A--H--M};
				\foreach \i/\g in {S/90,B/-90,C/0,A/180,M/-40,H/0,N/-140,G/-90}{\draw[fill=black](\i) circle (1.5pt) ($(\i)+(\g:4.5mm)$) node[scale=1]{$\i$};}
			\end{tikzpicture}
		\end{center}
		\begin{itemchoice}
			\itemch $M$ là trung điểm $BC$.\\
			Tam giác $ABC$ đều nên $AM\bot BC$ và $AM=\dfrac{a\sqrt{3}}{2}$.\\
			Ta có $SA \bot \left(ABC\right)\Rightarrow SA\bot AM$.\\
			Vậy khoảng cách giữa hai đường thẳng $SA$ và $BC$ là $AM=\dfrac{a\sqrt{3}}{2}$.\\
			Vậy phát biểu trên là sai.
			\itemch $N$ là trung điểm $AB$.\\
			Tam giác $ABC$ đều nên $CN\bot AB$ và $CN=\dfrac{a\sqrt{3}}{2}$\\
			$SA \bot \left(ABC\right)\Rightarrow SA\bot CN$.\\
			Vậy $CN\bot \left(SAB\right)$.\\
			Suy ra $\mathrm{d}\left(C,\left(SAB\right) \right)=CN=\dfrac{a\sqrt{3}}{2}$.\\
			Vậy phát biểu trên là đúng.
			\itemch Kẻ $AH\bot AM$ tại $H$.\\
			Ta có $BC\bot AM$\\
			$SA\bot \left(ABC\right)\Rightarrow SA\bot BC$.\\
			Ta có $BC\bot \left(SAM\right)\Rightarrow BC\bot AH$\\
			Vậy $AH\bot \left(SBC\right)$\\
			Suy ra $\mathrm{d}\left(A,\left(SAB\right) \right)=AH$\\
			Tam giác $SAM$ vuông tại $A$ có $\dfrac{1}{AH^2}=\dfrac{1}{SA^2}+\dfrac{1}{AM^2}=\dfrac{1}{a^2}+\dfrac{1}{{{\left(\dfrac{a\sqrt{3}}{2} \right)}^2}}=\dfrac{7}{3a^2}$\\
			Suy ra $AH^2=\dfrac{3a^2}{7}\Rightarrow AH=\dfrac{a\sqrt{21}}{7}$\\
			Vậy phát biểu trên là sai.
			\itemch Ta có $\dfrac{GM}{AM}=\dfrac{\mathrm{d}\left(G,\left(SBC\right) \right)}{\mathrm{d}\left(A,\left(SBC\right) \right)}=\dfrac{1}{3}$\\
			$\Rightarrow \mathrm{d}\left(G,\left(SBC\right) \right)=\dfrac{1}{3} \mathrm{d}\left(A,\left(SBC\right) \right)=\dfrac{1}{3}\cdot\dfrac{a\sqrt{21}}{7}=\dfrac{a\sqrt{21}}{21}$\\
			Vậy phát biểu trên là đúng.
	\end{itemchoice}}
\end{ex}
%%==============EX_4============%%%
\begin{ex}%[1H8V5-4]
	Cho hình lập phương $ABCD.A'B'C'D'$ có cạnh bằng $2$. Gọi $O$, $O'$ là tâm của hai đáy $ABCD$ và $A'B'C'D'$, $I$ là giao điểm của $BO'$ và $DB'$.
	\choiceTF
	{Khoảng cách giữa hai mặt $\left(ABCD\right)$ và $\left(A'B'C'D'\right)$ bằng $2$}
	{Khoảng cách từ $C$ đến $\left(BB'D'D\right)$ bằng đoạn $CO$}
	{Khoảng cách giữa hai đường thẳng $AD'$ và $DC'$ lớn hơn $2$}
	{Khoảng cách từ $B$ đến $\left(CDI\right)$ bằng đoạn $2\sqrt{2}$}
	\loigiai{
		\begin{center}
			\begin{tikzpicture}[line join=round,line cap=round,>=stealth,font=\footnotesize,scale=.8]
				\def\tls{.4}
				\def\r{7}
				\path
				(0,0) coordinate (A)
				++(0:\r) coordinate (D)
				++(-130:3) coordinate (C)
				($(A)+(C)-(D)$) coordinate (B)
				(intersection of A--C and B--D) coordinate (O)
				\foreach \x in {A,B,C,D,O}{(\x)++(90:\r) coordinate (\x\rq)}
				;
				\coordinate (I) at (intersection of O\rq--B and B\rq--D);
				\coordinate (H) at (intersection of B\rq--C and C\rq--B);
				\draw [dashed]
				(D)--(A)--(B)--(D)
				(C)--(A)--(A\rq)  (B\rq)--(D) (B)--(O\rq)
				(C)--(I) (D\rq)--(A)
				;
				\draw
				(D)--(C)--(B)
				(A\rq)--(B\rq)--(C\rq)--(D\rq)--cycle
				(D)--(D\rq) (B)--(B\rq) (C)--(C\rq) (B\rq)--(D\rq)
				(C\rq)--(D) (B\rq)--(C) (C\rq)--(B) 
				;
				\foreach \i/\g in {A/170,B/-150,C/-45,D/0,A\rq/120,D\rq/45,B\rq/180,C\rq/0,O/-90,O\rq/90,I/110,H/-90}{\draw[fill=black](\i) circle (1.5pt) ($(\i)+(\g:4mm)$) node[scale=1]{$\i$};}
			\end{tikzpicture}
		\end{center}
		\begin{itemchoice}
			\itemch Đúng.\\
			Vì ta có $\heva{&AA'\bot \left(ABCD\right) \\&AA'\bot \left(A'B'C'D' \right)}
			\Rightarrow \mathrm{d}\left(\left(ABCD\right),\left(A'B'C'D' \right) \right)=AA'=2$.
			\itemch Đúng. Vì ta có $\heva{&CO\bot BD\\&CO\bot BB'} \Rightarrow CO\bot \left(BB'D'D\right)\Rightarrow \mathrm{d}\left(C,\left(BB'D'D\right) \right)=CO$.
			\itemch Sai.\\
			Ta có
			$AD' \parallel BC' \Rightarrow AD' \parallel \left(BDC'\right)\Rightarrow \mathrm{d}\left(AD',DC' \right)=\mathrm{d}\left(AD',\left(BDC' \right) \right)=\mathrm{d}\left(A,\left(BDC' \right) \right)=\mathrm{d}\left(C,\left(BDC' \right) \right)$.
			Mà $\dfrac{1}{\mathrm{d}^2\left(C,\left(BDC' \right) \right)}=\dfrac{1}{CB^2}+\dfrac{1}{CD^2}+\dfrac{1}{C{{C'}^2}}=\dfrac{1}{2^2}+\dfrac{1}{2^2}+\dfrac{1}{2^2}=\dfrac{3}{4}\Rightarrow \mathrm{d}\left(C,\left(BDC' \right) \right)=\dfrac{2\sqrt{3}}{3}$.
			Vậy $\mathrm{d}\left(AD',DC' \right)=\dfrac{2\sqrt{3}}{3} < 2$.
			\itemch Sai.\\
			Ta có $I\in DB'\Rightarrow \left(CDI\right)\equiv \left(CDB' \right)\Rightarrow \mathrm{d}\left(B,\left(CDI\right) \right)=\mathrm{d} \left(B,\left(CDB\right) \right)$.
			Gọi $BC'\cap CB'=H\Rightarrow \heva{&BH\bot CB' \\
				&BH\bot CD}\Rightarrow BH\bot \left(CDB' \right)\Rightarrow \mathrm{d}\left(B,\left(CDB' \right) \right)=BH$.
			Vậy $\mathrm{d}\left(B,\left(CDI\right) \right)=BH=\dfrac{BC'}{2}=\dfrac{2\sqrt{2}}{2}=\sqrt{2}$.
	\end{itemchoice}}
\end{ex}
\Closesolutionfile{ans}

\ind{PHẦN III.} \inden{Câu trắc nghiệm trả lời ngắn.}\\
\Opensolutionfile{ans}[ans/1H7-B5-3]
\begin{ex}%[1H8V5-4]
	Cho tứ diện đều $ABCD$ có cạnh $2$. Tìm khoảng cách giữa hai đường thẳng chéo nhau $AB$, $CD$. (Kết quả làm tròn đến hàng phần trăm)\\
	\shortans[oly]{$1{,}41$}
	\loigiai{
		\begin{center}
			\begin{tikzpicture}[line join=round,line cap=round,>=stealth,font=\footnotesize,scale=1]
				\coordinate (A) at (2,4);
				\coordinate (B) at (-2,0);
				\coordinate (C) at (1,-2);
				\coordinate (D) at (4,0);
				\coordinate (I) at ($(A)!.5!(B)$);
				\coordinate (J) at ($(C)!.5!(D)$);
				\draw (C)--(I) (J)--(A)--(B)--(C) (A)--(C) (A)--(D) (C)--(D);
				\draw[dashed] (I)--(J)--(B)--(D)--cycle;
				\pic[draw,angle radius=2mm] {right angle = C--I--B};
				\pic[draw,angle radius=2mm] {right angle = D--I--A};
				\pic[draw,angle radius=2mm] {right angle = A--J--D};
				\pic[draw,angle radius=2mm] {right angle = B--J--C};
				\foreach \i/\g in {B/180,A/90,D/0,C/-90,I/130,J/-30}{\draw[fill=black](\i) circle (1.5pt) ($(\i)+(\g:3mm)$) node[scale=1]{$\i$};}
			\end{tikzpicture}
		\end{center}
		Gọi $I$, $J$ theo thứ tự là trung điểm của $AB$, $CD$.\\
		Các tam giác $ABC$, $ABD$ đều có $I$ là trung điểm $AB$ nên\\
		$\heva{&AB\bot CI\\&AB \bot DI} \Rightarrow AB \bot (ICD)$, mà $IJ\subset (ICD)\Rightarrow AB\bot IJ$. (1)\\
		Tương tự, các tam giác $ACD$, $BCD$ đều có $J$ là trung điểm $CD$ nên\\
		$\heva{&CD\bot AJ\\&CD\bot BJ} \Rightarrow CD \bot (ABJ)$,\\
		mà $IJ\subset (JAB)\Rightarrow CD\bot IJ (2)$\\
		Từ (1) và (2) suy ra $IJ$ là đoạn vuông góc chung của hai đường thẳng $AB$, $CD$.\\
		Ta có: $CI=\dfrac{2\cdot \sqrt{3}}{2}=\sqrt{3};IJ=\sqrt{CI^2-CJ^2}=\sqrt{3-1}=\sqrt{2}\approx 1{,}41$.}
\end{ex}

\begin{ex}%[1H8V5-3]
	Cho hình chóp $S.ABCD$, đáy $ABCD$ là hình thoi cạnh $2$. Tam giác $ABC$ đều, hình chiếu vuông góc $H$ của đỉnh $S$ trên mặt phẳng $ABCD$ trùng với trọng tâm tam giác $ABC$. Đường thẳng $SD$ hợp với mặt phẳng đáy một góc $30^\circ$. Tính khoảng cách từ $B$ đến mặt phẳng $\left(SCD\right)$. (làm tròn đến hàng phần trăm)\\
	\shortans[oly]{$1{,}31$}
	\loigiai{
		\begin{center}
			\begin{tikzpicture}[line join=round,line cap=round,>=stealth,font=\footnotesize,scale=.8]
				\coordinate (A) at (0,0);
				\coordinate (B) at (-2,-1.5);
				\coordinate (D) at (6,0);
				\coordinate (C) at ($(B)+(D)-(A)$);
				\coordinate (T) at ($(A)!.5!(B)$);
				\coordinate (O) at (intersection of A--C and B--D);
				\coordinate (H) at (intersection of T--C and B--D);
				\coordinate (S) at ($(H)+(0,6)$);
				\coordinate (K) at ($(S)!(H)!(C)$);
				\draw(S)--(B) (S)--(C) (S)--(D) (B)--(C)--(D);
				\pic[draw,angle radius=2mm] {right angle = H--K--S};
				\draw[dashed] (C)--(H)--(S)--(A) (A)--(B) (A)--(D) (B)--(D) (A)--(C) (H)--(K);
				\foreach \i/\g in {S/90,A/140,B/180,C/-45,D/0,O/50,H/150,K/70}{\draw[fill=black](\i) circle (1.5pt) ($(\i)+(\g:4mm)$) node[scale=1]{$\i$};}
			\end{tikzpicture}
		\end{center}
		Gọi $H$ là trọng tâm tam giác đều $ABC \Rightarrow SH \bot \left(ABCD\right)$\\
		$\Rightarrow HD$ là hình chiếu của $SD$ lên $\left(ABCD\right)$.\\
		$\Rightarrow \widehat{\left(SD;\left(ABCD\right) \right)}=\widehat{\left(SD;HD\right)}=\widehat{SDH}={{30}^\circ}$\\
		Ta có: $\dfrac{BH}{BO}=\dfrac{2}{3} \Rightarrow \dfrac{BH}{BD}=\dfrac{1}{3} \Rightarrow \dfrac{HD}{BD}=\dfrac{2}{3}$.\\
		Mà $BH \cap \left(SCD\right)= D \Rightarrow \dfrac{\mathrm{d}\left(B;\left(ABCD\right) \right)}{\mathrm{d}\left(H;\left(SCD\right) \right)}=\dfrac{BD}{HD}=\dfrac{3}{2}$\\
		$\Rightarrow \mathrm{d}\left(B;\left(SCD\right) \right)=\dfrac{3}{2} \mathrm{d}\left(H;\left(SCD\right) \right)$.\\
		Ta có $H$ là trọng tâm tam giác đều $ABC$ nên $HC\bot AB$.\\
		Mà $AB \parallel CD$ nên $HC\bot CD$.\\
		Trong $\left(SHC\right)$ kẻ $HK\bot SC$.\\
		Ta có $\heva{&CD\bot CH\\& CD\bot SH} \Rightarrow CD \bot \left(SHC\right)\Rightarrow CD\bot HK$.\\
		$\heva{& HK\bot SC\\& HK\bot CD} \Rightarrow HK\bot \left(SCD\right) \Rightarrow \mathrm{d}\left(H;\left(SCD\right) \right)=HK$.\\
		Vì tam giác $ABC$ đều cạnh bằng $2$ nên $HC=\dfrac{2}{3}\cdot\dfrac{2\sqrt{3}}{2}=\dfrac{2\sqrt{3}}{3}$.\\
		Ta có: $\dfrac{HD}{BD}=\dfrac{2}{3}$, mà $BD=2BO=2\cdot\dfrac{2\sqrt{3}}{2}=2\sqrt{3}\Rightarrow HD=\dfrac{2}{3}BD=\dfrac{2\cdot 2 \sqrt{3}}{3}=\dfrac{4\sqrt{3}}{3}$.\\
		Xét tam giác vuông $SHD$ ta có:\\ $SH=HD\cdot\tan{30^\circ}=\dfrac{4\sqrt{3}}{3}\cdot\dfrac{\sqrt{3}}{3}=\dfrac{4}{3}$.\\
		Áp dụng hệ thức lượng trong tam giác vuông $SHC$ ta có:\\
		$\dfrac{1}{HK^2}=\dfrac{1}{SH^2}+\dfrac{1}{HC^2}=\dfrac{1}{\dfrac{{{4\cdot2}^2}}{9}}+\dfrac{1}{\dfrac{2^2}{3}}=\dfrac{21}{16}$.\\
		$\Rightarrow HK=\dfrac{4\sqrt{21}}{21}$.\\
		Vậy $\mathrm{d}\left(B;\left(SCD\right) \right)=\dfrac{3}{2}HK=\dfrac{6\sqrt{21}}{21}\approx 1{,}31$}
\end{ex}
%%%==============EX_4============%%%
\begin{ex}%[1H8V5-4]
	Cho hình chóp $S.ABCD$ có đáy $ABCD$ là hình chữ nhật với $AB=a$, $AD=2a$. Hình chiếu vuông góc của $S$ trên mặt phẳng đáy là trung điểm $H$ của $AD$, góc giữa $SB$ và mặt phẳng đáy $\left(ABCD\right)$ là $45^\circ$. Khoảng cách giữa hai đường thẳng $SD$ và $BH$ theo $a$ được kết quả $ma$. Khi đó giá trị $\dfrac{3}{5}\, \mathrm{m}^2$ là bao nhiêu? (làm tròn đến hàng phần trăm)\\
	\shortans[oly]{$0{,}2$}
	\loigiai{
		\begin{center}
			\begin{tikzpicture}[line join=round,line cap=round,>=stealth,font=\footnotesize,scale=.8]
				\coordinate (A) at (0,0);
				\coordinate (B) at (-3,-2);
				\coordinate (D) at (7,0);
				\coordinate (C) at ($(B)+(D)-(A)$);
				\coordinate (T) at ($(A)!.5!(B)$);
				\coordinate (H) at ($(A)!.5!(D)$);
				\coordinate (E) at ($(B)!.5!(C)$);
				\coordinate (S) at ($(H)+(0,6)$);
				\coordinate (K) at ($(D)!.3!(E)$);
				\coordinate (I) at ($(S)!.8!(K)$);
				\draw(S)--(B) (S)--(C) (S)--(D) (B)--(C)--(D);
				\pic[draw,angle radius=2mm] {right angle = H--K--E}
				pic[draw,angle radius=2mm] {right angle = E--H--D}
				pic[draw,angle radius=2mm] {right angle = S--H--D}
				pic[draw,angle radius=2mm] {right angle = H--I--S}
				;
				\draw[dashed] (I)--(H)--(S)--(A) (A)--(B)--(H) (A)--(D)--(E)--(H)--(K)--(S);
				\foreach \i/\g in {S/90,A/170,B/180,C/-40,D/0,H/130,K/-80,E/-90,I/30}{\draw[fill=black](\i) circle (1.5pt) ($(\i)+(\g:4mm)$) node [scale=1] {$\i$};}
			\end{tikzpicture}
		\end{center}
		Ta có $SH\bot \left(ABCD\right)\Rightarrow$ góc giữa đường thẳng $SB$ và mặt phẳng đáy $\left(ABCD\right)$ là $\widehat{SBH}=45^\circ$.\\
		Suy ra $\Delta SBH$ vuông cân tại $H\Rightarrow$ $SH=BH=\sqrt{HA^2+AB^2}=a\sqrt{2}$.\\
		Gọi $E$ là trung điểm $CB$.\\
		Ta có $BH \parallel DE\Rightarrow \mathrm{d} \left(BH,SD\right)=\mathrm{d}\left(BH,\left(SDE\right) \right)= \mathrm{d}\left(H,\left(SDE\right) \right)$.\\
		Kẻ $HK\bot DE$, $HI\bot SK$.\\
		Ta có $DE\bot \left(SHK\right)\Rightarrow DE\bot HI$.\\
		Suy ra $HI\bot \left(SDE\right)$.\\
		Vậy $\mathrm{d}\left(BH,SD\right)= \mathrm{d}\left(H,\left(SDE\right) \right)=HI$.\\
		Trong $\Delta DHE$ vuông tại $H$ ta có $HK\cdot DE=DH\cdot HE\Leftrightarrow HK=\dfrac{DH \cdot HE}{DE}=\dfrac{a\cdot a}{a\sqrt{2}}=\dfrac{a\sqrt{2}}{2}$.
		Trong $\Delta SHK$ vuông tại $H$ ta có\\
		$\dfrac{1}{HI^2}=\dfrac{1}{SH^2}+\dfrac{1}{HK^2}\Leftrightarrow HI=\dfrac{SH\cdot HK}{\sqrt{SH^2+HK^2}}=\dfrac{a\cdot\dfrac{a\sqrt{2}}{2}}{\sqrt{a^2+\dfrac{2a^2}{4}}}=\dfrac{a}{\sqrt{3}}$.\\
		Vậy $\mathrm{d}\left(SD,BH\right)=\dfrac{a}{\sqrt{3}}$. Suy ra $m=\dfrac{1}{\sqrt{3}}$ nên $\dfrac{3}{5}\,\mathrm{m}^2=0{,}20$.}
\end{ex}
%%%==============EX_5============%%%
\begin{ex}%[1H8V5-4]
	Cho tứ diện $ABCD$, $AB=AD=2a$, $CD=2a\sqrt{2}$ và $\widehat{ABC}=\widehat{DAB}=90^\circ$. Góc giữa hai đường thẳng $AD$ và $BC$ bằng $45^\circ$. Khoảng cách giữa hai đường thẳng $AC$ và $BD$ bằng $\dfrac{a\sqrt{2k}}{k}$. Tìm $k$.\\
	\shortans[oly]{$3$}
	\loigiai{
		\begin{center}
			\begin{tikzpicture}[line join=round,line cap=round,>=stealth,font=\footnotesize,scale=.8]
				\coordinate (H) at (0,0);
				\coordinate (A) at (-2,-3);
				\coordinate (C) at (7,3);
				\coordinate (B) at (3,-3);
				\coordinate (M) at ($(B)!.5!(C)$);
				\coordinate (D) at ($(H)+(0,5)$);
				\coordinate (E) at ($(D)!.4!(B)$);
				\coordinate (I) at (intersection of A--C and H--B);
				\coordinate (U) at ($(E)+(I)-(H)$);
				\coordinate (F) at (intersection of D--B and I--U);
				\draw(D)--(A) (D)--(B)--(C)--(D)--(M) (A)--(B)--(M);
				\path (D)--(A) node [midway,left] {$2a$}
				(D)--(C) node [midway,above] {$2a\sqrt{2}$}
				(A)--(B) node [midway,below] {$2a$}
				;
				\draw[dashed](D)--(H)--(B) (H)--(A)--(C)--(H)--(M) (H)--(E) (I)--(F);
				\foreach \i/\g in {D/90,H/150,A/180,B/-90,M/0,C/0,I/-90,E/30,F/30}{\draw[fill=black](\i) circle (1.5pt) ($(\i)+(\g:3mm)$) node[scale=1]{$\i$};}
			\end{tikzpicture}
		\end{center}
		Gọi $H$ là hình chiếu vuông góc của $D$ lên mặt phẳng $\left(ABC\right)$\\
		Ta có $\heva{& DH\bot AB\\& DA\bot AB} \Rightarrow AB\bot AH$.\\
		Mà $BC\bot AB\Rightarrow AH\parallel BC\Rightarrow \left(AD,BC\right)=\left(AD,AH\right)=\widehat{DAH}=45^\circ$.\\
		Vì $\Delta DAH$ vuông tại $H$ $\Rightarrow AH=DH=a\sqrt{2}$.\\
		$\Delta DAB$ vuông tại $A\Rightarrow BD=\sqrt{AB^2+AD^2}=2a\sqrt{2}$.\\
		Suy ra $DB=DC=2a\sqrt{2}\Rightarrow \Delta DBC$ cân tại $D$.\\
		Gọi $M$ là trung điểm của $BC$, ta có \\
		$\heva{& DM\bot BC\\& DH\bot BC} \Rightarrow BC\bot HM\Rightarrow AHMB$ là hình chữ nhật.\\
		$\Rightarrow AH=BM=\dfrac{1}{2}BC$.\\
		Trong hình thang $ABCH$ $\left(AH\parallel BC\right)$, gọi $I=AC\cap BH$.\\
		Theo định lí Thales ta có $\dfrac{IH}{IB}=\dfrac{AH}{BC}=\dfrac{1}{2}\Rightarrow IH=\dfrac{1}{3}HB$.\\
		Ta lại có $\Delta AHB$ vuông tại $A$, $HB=\sqrt{AH^2+AB^2}=a\sqrt{6}$.\\
		Có $HI\cdot HB=\dfrac{1}{3}HB^2=2a^2=AH^2\Rightarrow AI\bot HB$ hay $HB\bot AC$.\\
		Ta có $\heva{& BH\bot AC\\& DH\bot AC}\Rightarrow AC\bot BD$.\\
		Trong $\Delta DHB$ dựng $HE\bot DB$, $IF\parallel HE\Rightarrow IF\bot BD\Rightarrow \mathrm{d}\left(AC,BD\right)=IF$\\
		$\Delta DHB$ vuông tại $H\Rightarrow HE=\dfrac{HB\cdot HD}{\sqrt{HB^2+HD^2}}=\dfrac{a\sqrt{6}}{2}$.\\
		Trong $\Delta EHB$ có\\
		$IF\parallel HE\Rightarrow \dfrac{IF}{HE}=\dfrac{BI}{BH}=\dfrac{2}{3}\Rightarrow IF=\dfrac{2}{3}HE\Rightarrow \mathrm{d}\left(AC,BD\right)=\dfrac{2}{3}HE=\dfrac{a\sqrt{6}}{3}$.\\
		Vậy $k=3$.}
\end{ex}
%%%==============EX_6============%%%
\begin{ex}%[1H8V5-3]
	Cho hình chóp $S.ABCD$ có đáy $ABCD$ là hình thang vuông tại $A$ và $B$ với $AB=2$, $BC=\dfrac{3}{2}$, $AD=3$. Hình chiếu vuông góc của $S$ lên mặt phẳng $\left(ABCD\right)$ là trung điểm $H$ của $BD$. Biết góc giữa mặt phẳng $\left(SCD\right)$ và mặt phẳng $\left(ABCD\right)$ bằng $60^\circ$. Tính khoảng cách từ điểm $C$ đến mặt phẳng $\left(SBD\right)$ (làm tròn đến hàng phần trăm)\\
	\shortans[oly]{$0{,}83$}
	\loigiai{
		\begin{center}
			\begin{tikzpicture}[line join=round,line cap=round,>=stealth,font=\footnotesize,scale=1]
				\coordinate (A) at (0,0);
				\coordinate (B) at (1.5,-3);
				\coordinate (D) at (7,0);
				\coordinate (C) at ($(B)+.5*(D)-(A)$);
				\coordinate (M) at ($(A)!.5!(D)$);
				\coordinate (H) at ($(B)!.5!(D)$);
				\coordinate (S) at ($(H)+(0,7)$);
				\draw (S)--(A)--(B)--(S)--(C) (S)--(D) (B)--(C)--(D);
				\draw[dashed] (A)--(D)--(B) (M)--(C) (A)--(H)--(S);
				\pic[draw,angle radius=2mm] {right angle = S--H--D} pic[draw,angle radius=2mm] {right angle = S--H--A};
				\path (B)--(H) node[midway,sloped]{$||$};
				\path (D)--(H) node[midway,sloped]{$||$};
				\foreach \i/\g in {S/90,A/180,B/-115,C/-45,D/0,M/90,H/-110}{\draw[fill=black](\i) circle (1.5pt) ($(\i)+(\g:4mm)$) node[scale=1]{$\i$};}
			\end{tikzpicture}
			\begin{tikzpicture}[line join=round,line cap=round,>=stealth,font=\footnotesize,scale=1]
				\coordinate (A) at (0,0);
				\coordinate (B) at (0,-4);
				\coordinate (D) at (6,0);
				\coordinate (C) at ($(B)+.5*(D)-(A)$);
				\coordinate (M) at ($(A)!.5!(D)$);
				\coordinate (H) at ($(B)!.5!(D)$);
				\coordinate (I) at ($(B)!(C)!(D)$);
				\coordinate (K) at ($(A)!(B)!(H)$);
				\draw (H)--(A)--(B)--(B)--(C)--(D)
				(A)--(D)--(B)--(K)
				(M)--(C)--(I)
				;
				\pic[draw,angle radius=2mm] {right angle = M--C--B}
				pic[draw,angle radius=2mm] {right angle = A--B--C}
				pic[draw,angle radius=2mm] {right angle = B--A--M}
				pic[draw,angle radius=2mm] {right angle = B--I--C}
				pic[draw,angle radius=2mm] {right angle = B--K--A};
				\foreach \i/\g in {A/150,B/-115,C/-45,D/0,M/90,H/-45,I/90,K/90}{\draw[fill=black](\i) circle (1.5pt) ($(\i)+(\g:4mm)$) node[scale=1]{$\i$};}
			\end{tikzpicture}
		\end{center}
		Ta có $SH \bot \left(ABCD\right) \Rightarrow \left(SBD\right) \bot \left(ABCD\right)$ theo giao tuyến $AH$.\\
		Trong mặt phẳng $\left(ABCD\right)$, kẻ $BK\bot AH$ tại $K$, ta có \\
		$BK\bot \left(SAH\right)\Rightarrow BK=\mathrm{d}\left(B,\left(SAH\right) \right)$.\\
		Tam giác $ABD$ vuông tại $A$ nên $BD^2=AB^2+AD^2=13\Rightarrow BD=\sqrt{13}\Rightarrow BH=\dfrac{\sqrt{13}}{2}$.\\
		Gọi $M$ là trung điểm của $AD$, ta có $\heva{& CM\parallel AB\\& HM\parallel AB}\Rightarrow$ ba điểm $C$, $H$, $M$ thẳng hàng.\\
		Ta có $HM=\dfrac{AB}{2}=1\Rightarrow CH=1$.\\
		Suy ra tam giác $BCH$ vuông tại $C$. Do đó, ta có
		$CI=\dfrac{CH \cdot CB}{BH}=\dfrac{3\sqrt{13}}{13}$.\\
		Vậy $\mathrm{d}\left(C,\left(SBD\right) \right)=CI=\dfrac{3\sqrt{13}}{13}\approx 0{,}83$.}
\end{ex}
\Closesolutionfile{ans}
